\documentclass[a4paper,11pt]{article}
\usepackage[utf8]{inputenc}
\usepackage{geometry}
\usepackage{xcolor}
\usepackage{listings}
\usepackage{hyperref}
\usepackage{graphicx}
\usepackage{amsmath}
\usepackage{amssymb}
\usepackage{tcolorbox}
\usepackage{float}
\usepackage{booktabs}

% Page geometry
\geometry{top=2.5cm, bottom=2.5cm, left=2.5cm, right=2.5cm}

% Custom colours
\definecolor{ibmblue}{RGB}{0, 45, 156}
\definecolor{ibmred}{RGB}{250, 75, 76}
\definecolor{purpleaccent}{RGB}{109, 40, 217}
\definecolor{amberaccent}{RGB}{245, 158, 11}
\definecolor{tealaccent}{RGB}{13, 148, 136}
\definecolor{shockorange}{RGB}{234, 88, 12}
\definecolor{erratumrose}{RGB}{190, 24, 93}
\definecolor{codegreen}{rgb}{0,0.6,0}
\definecolor{codegray}{rgb}{0.5,0.5,0.5}
\definecolor{codepurple}{rgb}{0.58,0,0.82}
\definecolor{backcolour}{rgb}{0.95,0.95,0.92}

% Python code style
\lstdefinestyle{mystyle}{
    backgroundcolor=\color{backcolour},
    commentstyle=\color{codegreen},
    keywordstyle=\color{magenta},
    numberstyle=\tiny\color{codegray},
    stringstyle=\color{codepurple},
    basicstyle=\ttfamily\footnotesize,
    breakatwhitespace=false,
    breaklines=true,
    captionpos=b,
    keepspaces=true,
    numbers=left,
    numbersep=5pt,
    showspaces=false,
    showstringspaces=false,
    showtabs=false,
    tabsize=2
}
\lstset{style=mystyle}

% Custom boxes (titles always in braced form)
\newtcolorbox{studentbox}{colback=blue!5!white,colframe=blue!75!black,title={Student Activity}}
\newtcolorbox{theorybox}{colback=gray!5!white,colframe=gray!50!black,title={Theory Recap}}
\newtcolorbox{warningbox}{colback=yellow!10!white,colframe=orange!80!black,title={Important}}
\newtcolorbox{simulatorbox}{colback=purple!5!white,colframe=purple!75!black,title={Simulator}}
\newtcolorbox{reflectionbox}{colback=teal!5!white,colframe=teal!75!black,title={Quick Check}}
\newtcolorbox{connectionbox}{colback=green!5!white,colframe=green!60!black,title={Connections to Previous Labs}}
\newtcolorbox{procedurebox}{colback=red!4!white,colframe=red!70!black,title={Procedure 1 vs Procedure 2}}
\newtcolorbox{questionbox}{colback=violet!4!white,colframe=violet!70!black,title={The Question}}
\newtcolorbox{shockbox}{colback=shockorange!8!white,colframe=shockorange,title={Counterintuitive --- and Real}}
\newtcolorbox{benchbox}{colback=teal!6!white,colframe=tealaccent,title={Open the Simulator First}}
\newtcolorbox{playbox}{colback=teal!6!white,colframe=tealaccent,title={Play First}}
\newtcolorbox{numberbox}{colback=green!4!white,colframe=green!50!black,title={Formula $\rightarrow$ Number $\rightarrow$ Bench}}
\newtcolorbox{industrybox}{colback=amberaccent!10!white,colframe=amberaccent!90!black,title={Real-World Anchor}}
\newtcolorbox{deepdivebox}{colback=cyan!4!white,colframe=cyan!60!black,title={Deep Dive}}
\newtcolorbox{erratumbox}{colback=erratumrose!5!white,colframe=erratumrose,title={Erratum}}
\newtcolorbox{creditsbox}{colback=amberaccent!8!white,colframe=amberaccent!80!black,title={About the Original Game}}

% Bra-ket shorthands (guarded: never clash with package-provided macros)
\providecommand{\ket}[1]{|#1\rangle}
\providecommand{\bra}[1]{\langle#1|}
\providecommand{\braket}[2]{\langle#1|#2\rangle}
\providecommand{\ketbra}[2]{|#1\rangle\langle#2|}

\title{\textbf{Laboratory Guide 8: QTris --- Quantum Mechanics as a Board Game}\\
\large A Verified Digital Adaptation of QTris}
\author{Course: Quantum Computing Foundation}
\date{}

\begin{document}

\maketitle

\section{Introduction and Objectives}

QTris plays tic-tac-toe on a $3 \times 3$ grid whose nine cells are \textbf{nine
qubits}. Every tile is a quantum state, every move card is a reversible operation, and
the endgame is a dice-resolved measurement that collapses the whole board to black and
white before a single point can be scored. Its rulebook makes a bold claim: \emph{there
is no fundamental difference between QTris and quantum mechanics.} That claim is
checkable --- and this lab checks it. Every probability in the digital adaptation you
will play is computed by an engine whose numbers were verified, state by state and table
by table, against the game's printed measurement tables. Where the print and the physics
agree (38 of the 40 printed rows, within a stated rounding of at most half a percentage
point), you play the print. Where they provably cannot both be right (two rows --- see
the erratum section), the physics wins, openly.

You will play first and formalize second. The tally counters, the printed dice tables
and an optional exact mode are all part of the instrument: the game is not an analogy
for quantum mechanics --- it \emph{is} a quantum system you operate by hand, the same
physics that runs Stern--Gerlach spin measurements and commercial quantum random-number
generators.

\vspace{0.3cm}
\noindent\textbf{Prerequisites:}
\begin{itemize}
    \item \textit{Fundamental Concepts of Quantum Technologies} (lecture series).
    \item Lab 5: \textit{Quantum Measurements} --- the Born rule, measurement
          Procedures 1 and 2, the states $\ket{0}, \ket{1}, \ket{+}, \ket{-}$ and the
          four Bell states.
    \item Lab 6: \textit{Quantum Key Distribution (BB84)} --- mutually unbiased bases
          (this lab's colour and orientation are exactly that pair).
    \item Lab 7: \textit{Quantum Optics} --- the invariance of $\ket{\Phi^+}$ when both
          members of a pair are rotated together, which this lab reuses and sharpens.
    \item Familiarity with Dirac notation: $\ket{\psi}$, $\bra{\phi}$, $\braket{\phi}{\psi}$.
\end{itemize}

\begin{connectionbox}
\textbf{The Lab 5--7 bridge.} The game's tiles are the Lab 5 states: white and black
are the computational basis $\ket{0}, \ket{1}$, the two half-and-half tiles are
$\ket{+}, \ket{-}$. Its colour and orientation questions are Lab 6's mutually unbiased
pair, $Z$ versus $X$. Its entangled pairs are Lab 7's photon pairs with the analyzers
replaced by dice, and the double-pink-square analysis in the erratum section is Lab 7's
``rotate both analyzers together'' invariance --- now shown to hold for exactly two of the
four Bell states.
\end{connectionbox}

\noindent\textbf{Mode:} Playable Board Game + Qiskit \quad\textbf{Credits:} 1 ECTS
\quad\textbf{Estimated time:} $\approx$ 8 h

\vspace{0.3cm}
\noindent\textbf{Learning Objectives:}
\begin{enumerate}
    \item \textbf{Translate every QTris component into quantum-mechanical language} ---
          tiles as states $\ket{W} = \ket{0}$, $\ket{B} = \ket{1}$, $\ket{R} = \ket{+}$,
          $\ket{L} = \ket{-}$; cards as unitaries; dice as the Born rule.
    \item \textbf{Predict measurement statistics} for single tiles and entangled pairs,
          decorated or not, and confirm them against both the printed d100 tables and
          the exact probabilities.
    \item \textbf{Explain why the pink-square gate is a single canonical rotation}
          $R(-60^\circ)$, and how one rotation produces the game's 25, 75, 93 and 7 ---
          including where the $\sqrt{3}$ comes from.
    \item \textbf{Distinguish Procedure 1 and Procedure 2} for a rotated-basis
          measurement: identical statistics, different post-measurement states --- on
          the game's own tiles.
    \item \textbf{Recognize state-dependent rotational invariance}: why two pink squares
          cancel exactly on two of the four entangled pair types and provably not on the
          other two --- and how that catches a printed table in an error its own authors
          missed.
    \item \textbf{Operate sequential incompatible measurements} (colour vs orientation)
          and connect them to Stern--Gerlach physics and spin qubits.
\end{enumerate}

\begin{creditsbox}
\textbf{QTris} (\emph{QTris e Meccanica Quantistica}) was created by \textbf{Alioscia
Hamma} (concept and physics) with the creative team \textbf{Immacolata De Simone},
\textbf{Michela Nazzaro} and \textbf{Michele Viscardi} (illustrations by M.~Viscardi),
produced by AH Quantum Lab and Communikart, and published online in 2024 by
\textbf{NQSTI} --- the National Quantum Science and Technology Institute
(\url{https://www.nqsti.it}), Italy's national quantum research partnership. The
original rulebook is in Italian and is the authoritative source for every rule and
printed table this adaptation implements. This lab is an independent educational
adaptation built inside the same national ecosystem, with respect and with every
deviation from the printed game declared explicitly (Section~\ref{sec:deviations}).
The rulebook's own philosophy is worth quoting in spirit: the rules are not a metaphor
for quantum mechanics --- the rules \emph{are} quantum mechanics. This adaptation takes
that philosophy literally enough to verify it. No artwork from the original is
reproduced here; only the printed numeric tables are transcribed, for the audit.
\end{creditsbox}

\noindent The original terms appear throughout the Italian rulebook; this guide uses
the English terms below.

\begin{table}[H]
\centering
\footnotesize
\begin{tabular}{p{3.6cm}p{4.6cm}p{6.6cm}}
\toprule
Italian (rulebook) & English (this guide) & Meaning \\
\midrule
tris & tris (three-in-a-row) & kept as the game's name element \\
griglia / casella & grid / cell & the $3\times 3$ board; cells numbered 1--9 \\
tessera & tile & the piece occupying a cell \\
bianconera & \emph{bianconera} --- the half-and-half tile & the white-and-black tile; kept as a loanword \\
orientata a destra / sinistra & right- / left-oriented & white half on the right / on the left \\
carta mossa & move card & a playable operation: $I, X, Y, Z, H, U, C_X$ \\
tessera quadrata rosa $U$ / rossa $H$ & pink $U$ square / red $H$ square & a marker placed under a tile \\
decorare & to decorate & placing a square marker under a tile \\
gattini di Schr\"odinger & Schr\"odinger's kittens & the rulebook's name for the bianconere \\
stato entangled & entangled state & one state on two cells (triangle tiles) \\
fase di preparazione / operazioni / misura & preparation / operations / measurement phase & the three phases of a game \\
mulligan & mulligan & one optional opening redraw (rulebook's own loanword) \\
Espansione & deck preset & Base, $+U$, $+C_X$, $+U{+}C_X$ \\
\bottomrule
\end{tabular}
\caption{Translation glossary.}
\label{tab:glossary}
\end{table}

\begin{simulatorbox}
The lab's interactive page (\texttt{lab8\_qtris.html}) hosts the \textbf{playable
game} built for this lab --- hot-seat, two players at one screen --- and it is the bench
of this guide. Its \textbf{Play QTris} tab offers the four official deck presets
(\textbf{Base}, \textbf{+U}, \textbf{+C$_X$} --- the didactic default ---
and \textbf{+U+C$_X$}), a \textbf{dice toggle} between \emph{printed-tables} mode
(d100 rolled against the published bands, the consulted row shown on every roll) and
\emph{exact} mode (full-precision Born sampling), a seeded random setup or the
rulebook's own example board, an opt-in probability preview, session tallies, and a
scoring overlay on the eight lines. Its \textbf{Sandbox} tab is a single cell with a
free card rack, colour measurement, and orientation measurement by Procedure~1 and by
Procedure~2. Every displayed probability flows through the same verified reference
engine as the formulas in this guide, and all runs are seeded: the same seed and the
same sequence of actions always reproduce the same game.

\medskip
\textbf{This guide is written play-first.} Each Part opens with a ``Play First'' box.
Do that two-minute run \emph{before} reading the theory: what you see at the board
poses exactly the question the section then answers. The companion Qiskit notebook
rebuilds the same states as quantum circuits and audits the printed tables in code:\\[4pt]
\url{https://colab.research.google.com/github/LMrilo/quantum-labs-2026/blob/main/lab8/lab8_qtris.ipynb}
\end{simulatorbox}

\newpage

\section{Part 1: The Board --- Nine Cells, Nine Qubits}

\begin{playbox}
Start a game in the default preset (\textbf{+C$_X$}) with the dice toggle on
\textbf{printed tables}. Play the ten turns any way you like, then run the measurement
sweep and watch the dice panel: every roll names the printed row it consults and the
band it lands in. Look at the session tally afterwards --- the fraction of half-and-half
tiles that came up white. Play a few games: it drifts toward \textbf{50\%}. Keep that
number in mind; Part~2 will do something to it that no coin allows.
\end{playbox}

\begin{questionbox}
Every cell is full from move one, and yet nobody knows who's winning. What kind of
board game can't tell you its own score until the dice say so?
\end{questionbox}

A game of QTris has three phases. In \textbf{preparation}, all nine cells are filled
with tiles. In the \textbf{operations phase}, the players alternate for ten turns (five
each), each turn drawing two cards and playing two --- cards act on any cell; there is
no cell ownership. In the \textbf{measurement phase}, the cells are visited in order
$1 \to 9$ and every unresolved tile is collapsed by a dice roll to plain white or
black. Only then is the board scored: one point per three-in-a-row of your colour over
the eight lines (three rows, three columns, two diagonals); most points wins.

The tiles are states of a two-level system --- a qubit per cell:

\begin{equation}
    \ket{W} = \ket{0}, \qquad \ket{B} = \ket{1}, \qquad
    \ket{R} = \ket{+} = \frac{\ket{0} + \ket{1}}{\sqrt 2}, \qquad
    \ket{L} = \ket{-} = \frac{\ket{0} - \ket{1}}{\sqrt 2}.
    \label{eq:tiles}
\end{equation}

White and black are the computational basis. The \emph{bianconera} --- half white,
half black, with an orientation --- is an equal superposition; drawn with the white half
on the \textbf{right} it is $\ket{+}$ (``right-oriented''), mirrored it is $\ket{-}$.
This orientation convention is the one under which every drawing and every passage of
the rulebook is self-consistent (one early caption is the sole exception --- erratum E3
in Section~\ref{sec:erratum}). The colour measurement asks a superposed tile a question
it does not have a definite answer to, and the Born rule replies with probabilities:

\begin{table}[H]
\centering
\begin{tabular}{llcc}
\toprule
tile & state & $P(W)$ & printed white band (d100) \\
\midrule
white & $\ket{0}$ & $1$ & 1--100 \\
black & $\ket{1}$ & $0$ & --- (no white band) \\
bianconera, right-oriented & $\ket{+}$ & $1/2$ & 1--50 \\
bianconera, left-oriented & $\ket{-}$ & $1/2$ & 1--50 \\
\bottomrule
\end{tabular}
\caption{The four undecorated tiles and their printed dice rows. A roll of 1--50 on a
bianconera is white, 51--100 black.}
\label{tab:tiles}
\end{table}

\begin{numberbox}
$P(W \mid R) = |\braket{0}{+}|^2 = \tfrac12$, i.e.\ $0.500$ --- and the printed dice
row for the bianconera is exactly 1--50 white, 51--100 black. At the board: measure a
few games with bianconere on them and watch the white tally converge to 50\%. Then
remember that this 50/50 is \emph{not} a coin: Part~2 un-flips it.
\end{numberbox}

\begin{theorybox}
\textbf{Global phase is discarded --- by the game, deliberately.} A quantum state and
the same state multiplied by a phase factor give identical probabilities for every
measurement, so the game identifies them. Two consequences you can check on the move
cards: the card $Z$ does nothing visible to a monochrome tile ($Z\ket{0} = \ket{0}$,
$Z\ket{1} = -\ket{1}$) and the card $X$ does nothing visible to a bianconera
($X\ket{\pm} = \pm\ket{\pm}$), so both are \emph{legal but invisible} plays --- the
digital game lets you make them and tells you so. The card $Y$ acts as ``$X$ and $Z$
at once'': $Y = iXZ$, and the $i$ is exactly what the game discards. Even the card $I$
--- do nothing --- is a legal tempo move. A move that provably does nothing is itself
teachable: it is what \emph{global phase} means.
\end{theorybox}

\begin{industrybox}
\textbf{A research institute's outreach instrument.} QTris is not a classroom toy
invented for this course. It is an instrument of Italy's national quantum initiative:
conceived by a physicist, produced by a design team, playtested by university
researchers and published by NQSTI --- the same ecosystem this course belongs to.
Game-based learning is an active research line in quantum education (see the survey by
Seskir et al.\ in the Resources), and QTris sits at its most literal end: the rules are
not a metaphor for the theory, they \emph{are} the theory, which is why an engine can
audit them.
\end{industrybox}

\section{Part 2: The Bianconera --- Why Orientation Matters}

\begin{playbox}
In any preset, find a right-oriented bianconera on the board (or make one: play $H$ on
a white tile). Now play $H$ on it and run the sweep. The tile needs no dice at all ---
it is white, with certainty. Do it with a left-oriented one: black, with certainty.
The rulebook itself recommends this tactic.
\end{playbox}

\begin{questionbox}
A half-white-half-black tile turns up white exactly half the time. So does a coin. Why
does the game bother printing an orientation on the tile --- what could the direction
of a 50/50 tile possibly change?
\end{questionbox}

The move cards answer it. The card $H$ connects the two families: it turns white into a
right-oriented bianconera and back, black into a left-oriented one and back. So play $H$
on a bianconera and the superposition \emph{closes}:

\begin{equation}
    H\ket{+} = \ket{0}, \qquad H\ket{-} = \ket{1}.
    \label{eq:hclose}
\end{equation}

\begin{shockbox}
\textbf{One card un-flips the coin.} No operation on a fair coin can make its next toss
certain --- a probability distribution has no handle to grab. But $H$ on a
right-oriented bianconera yields a \emph{certainly white} tile: 100\%, not 50. The
bianconera was never a distribution; it is a state with a definite \emph{orientation},
and the orientation is the handle. The rulebook calls these tiles Schr\"odinger's
kittens for exactly this reason: the superposition is a real, manipulable thing.
\end{shockbox}

The same logic exposes the phase. The card $Z$ flips orientation,
$Z\ket{\pm} = \ket{\mp}$, yet changes \emph{nothing} a colour measurement can see ---
both orientations measure 50/50. Invisible? Only until an $H$ converts the hidden
difference into opposite certainties: $H\ket{+} = \ket{0}$ (white for sure) while
$H\ket{-} = \ket{1}$ (black for sure).

\begin{shockbox}
\textbf{The $Z$ card changes nothing you can see --- until it changes everything.}
Phase is not bookkeeping. The game tracks it with tile orientation, and one $H$ card
cashes it out as opposite certain colours. This is the same relative phase that Lab 7's
interferometer converted into which-detector certainty.
\end{shockbox}

\begin{numberbox}
$H\ket{+} = \ket{0}$, so $P(W \mid H \text{ on } R) = 1$: the number is $1.000$ --- not
$0.500$. At the board: play $H$ on a bianconera, then measure. This is how you build a
guaranteed tris piece out of chance.
\end{numberbox}

\begin{deepdivebox}
\textbf{The single-tile card algebra.} The rulebook's operation map for single tiles is
the Pauli/Hadamard action on the dictionary of equation~\eqref{eq:tiles}, taken modulo
global phase: $X$ is the colour flip ($X\ket{0} = \ket{1}$; invisible on bianconere),
$Z$ is the orientation flip ($Z\ket{\pm} = \ket{\mp}$; invisible on monochrome tiles),
$Y$ does both at once ($Y = iXZ$, phase discarded), and $H$ is the two bridges
$W \leftrightarrow R$, $B \leftrightarrow L$. Every one of these cards is its own
inverse: $X^2 = Y^2 = Z^2 = H^2 = I$, so any card played twice on the same tile is a
round trip --- the rulebook's own unitarity note, made tactile. Playing $H$ twice on a
bianconera therefore returns the bianconera; playing $Z$ between the two $H$'s returns
the \emph{opposite} colour --- phase, converted twice.
\end{deepdivebox}

\begin{connectionbox}
Lab 7 converted a relative phase between two paths into a certain detector click at
the second beam splitter. Equation~\eqref{eq:hclose} is the same conversion with the
beam splitter replaced by an $H$ card and the two paths by the two colours: the
orientation of a bianconera is a phase, and $H$ is the device that reads it.
\end{connectionbox}

\section{Part 3: Entanglement --- One State on Two Cells}

\begin{playbox}
In preset \textbf{+C$_X$}, play the $C_X$ card with a bianconera as control and a
monochrome tile as target (the game asks you to declare which is which). The two cells
become a pair of triangles --- one state on two cells. Run the sweep, and watch the pair
counter in the session tally over several games: \textbf{100\% agreement}, while each
member's own colour tally sits at 50/50. Entangle cells 1 and 9 --- opposite corners ---
and check that distance changes nothing.
\end{playbox}

\begin{questionbox}
Two triangles on opposite corners of the board always come up the same colour, every
game, though each alone is a perfect coin toss. Where does the agreement live?
\end{questionbox}

The two-cell card $C_X$ (a controlled-NOT: the player declares which cell is control
and which is target) is the game's only source of correlation. Played with a bianconera
as control and a monochrome tile as target, it \emph{entangles} them into one of the
four Bell states, shown on the board as a pair of triangles:

\begin{equation}
    (R, W) \to \Phi^+, \qquad (R, B) \to \Psi^+, \qquad
    (L, W) \to \Phi^-, \qquad (L, B) \to \Psi^-
    \label{eq:bellrecipe}
\end{equation}

\noindent --- control orientation $\mapsto$ sign of the pair, target colour $\mapsto$
same-colour ($\Phi$) or opposite-colour ($\Psi$) type. The recipe is one line of algebra:
with the control written first,

\begin{equation}
    \ket{+}\ket{0} = \frac{\ket{00} + \ket{10}}{\sqrt 2}
    \;\xrightarrow{\;C_X\;}\;
    \frac{\ket{00} + \ket{11}}{\sqrt 2} = \ket{\Phi^+},
    \label{eq:cxderiv}
\end{equation}

\noindent and likewise $\ket{+}\ket{1} \to (\ket{01} + \ket{10})/\sqrt 2 = \ket{\Psi^+}$,
$\ket{-}\ket{0} \to (\ket{00} - \ket{11})/\sqrt 2 = \ket{\Phi^-}$,
$\ket{-}\ket{1} \to (\ket{01} - \ket{10})/\sqrt 2 = \ket{\Psi^-}$.

\medskip
\noindent\textbf{Reading order} (used everywhere in this guide, on the interactive page
and in the notebook): pair outcomes like $WB$ are written \textbf{big-endian} --- the
lower-numbered cell is the \textbf{leftmost} letter, exactly as $\ket{q_0 q_1}$ puts
qubit 0 leftmost. Amplitude vectors are listed in the order
$(WW, WB, BW, BB) = (\ket{00}, \ket{01}, \ket{10}, \ket{11})$. The companion notebook
shows how to convert the opposite ordering used by common circuit software.

A $\Phi^+$ pair measures $WW$ or $BB$ --- never mixed --- each with probability
$\tfrac12$: each member alone is perfectly random, the two together perfectly agreed.
The agreement does not live in either cell; it lives in the pair. Adjacency is
irrelevant --- cells 1 and 9 entangle as happily as neighbours. And the same card played
again on the pair \emph{disentangles} it back into the tiles it came from (the player
chooses which output tile lands on which cell): every operation in the game is
reversible; only the dice are not.

\begin{shockbox}
\textbf{Two cells that never touched agree forever.} One $C_X$ card creates perfect
correlation with individually perfect randomness, at any board distance. This is the
same resource measured in Lab 7's entangled-photon correlations, and the resource behind
quantum key distribution (Lab 6): the correlation is a property of the joint state,
stored in neither place.
\end{shockbox}

\begin{numberbox}
$P(\text{same colour} \mid \Phi^+) = |\braket{00}{\Phi^+}|^2 + |\braket{11}{\Phi^+}|^2
= \tfrac12 + \tfrac12 = 1$: agreement $1.000$, and the mixed outcomes have probability
exactly $0$ --- bands that simply do not exist on the printed row (1--50 $\to WW$,
51--100 $\to BB$, nothing else). At the board (preset +C$_X$): entangle two cells,
measure many games, and watch the pair counter pin agreement at 100\% while each
member's own tally sits at 50/50.
\end{numberbox}

\begin{deepdivebox}
\textbf{The marker calculus --- how decorations compose.} A pair member may carry
markers under its triangle: a red $H$ square (the $H$ card played on that member) and a
pink $U$ square (Part~4). Each member's stack is one of $\{I, H, U, UH\}$, read from the
tile outward, with $U$ always outermost: the $H$ card toggles the layer next to the
tile, the $U$ card toggles the top layer. Three identities make the printed operation
tables come out exactly:
\begin{itemize}
    \item \textbf{Both members marked $H$ normalize away.} $(H \otimes H)$ leaves
          $\ket{\Phi^+}$ and $\ket{\Psi^-}$ unchanged and \emph{swaps} the other two:
          $(H \otimes H)\ket{\Phi^-} = \tfrac{1}{\sqrt 2}(\ket{{+}{+}} - \ket{{-}{-}})
          = \tfrac{1}{\sqrt 2}(\ket{01} + \ket{10}) = \ket{\Psi^+}$. So when the second
          member receives an $H$ marker too, both markers come off and the pair is
          relabelled $\Phi^- \leftrightarrow \Psi^+$ --- the rulebook prints exactly this
          column.
    \item \textbf{Paulis commute through $H$ with $X \leftrightarrow Z$.} $HX = ZH$ and
          $HZ = XH$ (and $Y$ passes through up to phase). A $Z$ card on an $H$-marked
          member therefore acts as $X$ on the core --- which is why the rulebook needs
          \emph{two} $X$ columns and \emph{two} $Z$ columns for marked pairs (played on
          the marked vs the unmarked member) but only one $Y$ column.
    \item \textbf{Phase kickback.} $C_X$ with a bianconera control \emph{and} a
          bianconera target flips the \emph{control's} orientation exactly when the
          target is left-oriented: $C_X\,\ket{+}\ket{-} = \ket{-}\ket{-}$, while a
          right-oriented target is an eigenstate and the play is invisible. The X-basis
          CNOT runs backwards --- the case the printed tables do not cover, fixed by the
          physics.
\end{itemize}
\textbf{Counting the state space.} 8 single-cell states (4 tiles $\times$ with/without
pink square) and $4 \times 12 = 48$ entangled states (4 Bell cores $\times$ the
$4 \times 4 = 16$ stack pairs minus the 4 in which both members carry $H$, which
normalize away): 56 states in all, closed under every card. Small enough that every
claim in this guide was checked exhaustively, not sampled.
\end{deepdivebox}

\begin{theorybox}
\textbf{No signalling, printed.} Whatever markers sit under \emph{one} member of a pair,
the \emph{other} member alone is a perfect coin: both single-cell marginals of every
pair row are exactly $\tfrac12$, because local operations cannot change the reduced
state $I/2$ of a Bell pair (Lab 7's ``each photon alone is a perfect coin''). The
printed tables respect this \emph{even after rounding}: where the exact eighths
$12.5/37.5$ had to become integers, the authors rounded them in opposite directions
(12 and 38) so that every row still tiles 1--100 and both marginals still sum to
exactly 50. Table~\ref{tab:pairs} shows the result.
\end{theorybox}

\begin{connectionbox}
Lab 7's Part D showed photon pairs agreeing at matched analyzer angles while each side
alone stayed 50/50; Lab 6 used the same correlations as the raw material of a key. The
$C_X$ card is the pair source, the sweep is the pair of analyzers, and the printed
$\Phi^+$ row --- $WW$ or $BB$, nothing else --- is the matched-angle perfect correlation.
\end{connectionbox}

\section{Part 4: The Pink Square --- One Hidden Rotation \textnormal{\normalsize(advanced)}}

\begin{playbox}
Switch to preset \textbf{+U}. Play the $U$ card on a right-oriented bianconera --- a pink
square appears under it --- and run the sweep in printed-tables mode: the dice panel
consults a \textbf{1--93} row. Flip the toggle to \textbf{exact} mode and measure again:
the panel now serves $0.933013$ with all its digits. Decorate a white tile instead:
\textbf{1--25}.
\end{playbox}

\begin{questionbox}
One pink square turns your certain white tile into a 25\% underdog. What single
operation produces 25, 75, 93 \emph{and} 7 --- and why does the 93 come with a
$\sqrt{3}$ in it?
\end{questionbox}

The rulebook never writes the pink square's matrix --- it only prints its measurement
table: a decorated white tile measures white 25\%, decorated black 75\%, decorated
bianconere 93\% and 7\%. Those four rows pin the gate completely. Take the square to be
a canonical (real) rotation by an angle $\theta$,

\begin{equation}
    R(\theta) = \begin{pmatrix} \cos\theta & -\sin\theta \\ \sin\theta & \cos\theta \end{pmatrix},
    \qquad R(\theta)\ket{0} = \cos\theta\,\ket{0} + \sin\theta\,\ket{1}.
    \label{eq:rot}
\end{equation}

\noindent The decorated-white row demands $\cos^2\theta = \tfrac14$, hence
$\theta = \pm 60^\circ$; matching the 93 to the \emph{right}-oriented bianconera (and
the 7 to the left-oriented one) fixes the sense:

\begin{equation}
    U = R(-60^\circ) = \begin{pmatrix} 1/2 & \sqrt{3}/2 \\ -\sqrt{3}/2 & 1/2 \end{pmatrix}.
    \label{eq:ugate}
\end{equation}

\noindent Now $\ket{+} = R(45^\circ)\ket{0}$, and rotations compose by adding angles, so

\begin{equation}
    U\ket{+} = R(-60^\circ)\,R(45^\circ)\ket{0} = R(-15^\circ)\ket{0},
    \label{eq:u15}
\end{equation}
\begin{equation}
    P(W \mid U \text{ on } R) = \cos^2 15^\circ = \frac{1 + \cos 30^\circ}{2}
    = \frac{2 + \sqrt{3}}{4} \approx 0.9330.
    \label{eq:cos15}
\end{equation}

The game's mysterious 93 is $\cos^2$ of fifteen degrees --- an irrational probability,
printed to two digits. Table~\ref{tab:single} lists all eight single-cell rows against
the print.

\begin{table}[H]
\centering
\small
\begin{tabular}{llccc}
\toprule
tile & state & exact $P(W)$ & printed white band & gap \\
\midrule
white & $\ket{0}$ & $1$ & 1--100 & 0 \\
black & $\ket{1}$ & $0$ & --- & 0 \\
bianconera right & $\ket{+}$ & $1/2$ & 1--50 & 0 \\
bianconera left & $\ket{-}$ & $1/2$ & 1--50 & 0 \\
white $+\,U$ & $U\ket{0}$ & $1/4$ & 1--25 & 0 \\
black $+\,U$ & $U\ket{1}$ & $3/4$ & 1--75 & 0 \\
bianconera right $+\,U$ & $U\ket{+} = R(-15^\circ)\ket{0}$ & $(2+\sqrt 3)/4 \approx 0.9330$ & 1--93 & $0.0030$ \\
bianconera left $+\,U$ & $U\ket{-}$ & $(2-\sqrt 3)/4 \approx 0.0670$ & 1--7 & $0.0030$ \\
\bottomrule
\end{tabular}
\caption{The eight single-cell states and their printed dice rows. The two irrational
rows are rounded by the print to two digits; the gap is $0.0030$.}
\label{tab:single}
\end{table}

Unlike $X$, $Y$, $Z$ and $H$, this rotation does not map the simple tile family onto
itself ($U\ket{0}$ is none of the four tiles) --- which is exactly why the pink square
must \emph{stay on the board} as a marker instead of being absorbed into a relabelled
tile, and why the +U presets are the advanced tier: one little square carries the game
beyond the flip-and-mix world into genuinely new probabilities. In the language of Lab 5,
$U$ is not a Clifford operation: it does not send Pauli operators to Pauli operators.

\begin{numberbox}
$P(W \mid U \text{ on } R) = \cos^2 15^\circ = \tfrac{2+\sqrt 3}{4}$. As a float64
number this is $0.9330127018922193$ (exact mode displays $0.933013$); the print's 93
rounds the physics by $0.0030$, within its stated half-point tolerance. At the board
(preset +U): decorate a bianconera, roll against the printed 93/7 row, then flip to exact
mode and see the value served with all its digits.
\end{numberbox}

\begin{deepdivebox}
\textbf{Where the stack order is pinned, and where the eighths come from.} Two facts
about decorated \emph{pairs} follow from equation~\eqref{eq:ugate}.
\begin{itemize}
    \item \textbf{$U$ outermost, not innermost.} A member carrying both a red and a pink
          square could mean $UH$ (apply $H$, then $U$) or $HU$. The two differ:
          $|\bra{0}UH\ket{0}|^2 = |\bra{0}U\ket{+}|^2 = \tfrac{2+\sqrt 3}{4}$ while
          $|\bra{0}HU\ket{0}|^2 = \tfrac{2-\sqrt 3}{4}$. On a $\Phi^+$ pair the first
          gives the outcome probabilities $(q_+, q_-, q_-, q_+)$ with
          $q_\pm = \tfrac{2 \pm \sqrt 3}{8} \approx 0.4665,\ 0.0335$ --- the printed
          \textbf{47/3/3/47} row --- and the second gives $(q_-, q_+, q_+, q_-)$, i.e.\
          3/47/47/3. The print decides: the pink square sits outermost.
    \item \textbf{The eighth family.} With a pink square on member 1 only,
          $(U \otimes I)\ket{\Phi^+} = \tfrac{1}{\sqrt 2}\bigl(U\ket{0}\ket{0} +
          U\ket{1}\ket{1}\bigr)$, and $U\ket{0} = (\tfrac12, -\tfrac{\sqrt 3}{2})$,
          $U\ket{1} = (\tfrac{\sqrt 3}{2}, \tfrac12)$ give the amplitude vector
          $\tfrac{1}{2\sqrt 2}\,(1, \sqrt 3, -\sqrt 3, 1)$ over $(WW, WB, BW, BB)$,
          hence the probabilities $(\tfrac18, \tfrac38, \tfrac38, \tfrac18)$ --- the
          rows the rulebook resolves exactly with an eight-sided die.
    \item \textbf{Mirror symmetry.} Swapping the two members' stacks swaps the $WB$ and
          $BW$ columns and nothing else, which is why the rulebook may tabulate
          marker-on-member-1 configurations without loss.
\end{itemize}
\end{deepdivebox}

\begin{table}[H]
\centering
\small
\begin{tabular}{llcc}
\toprule
pattern & exact $(WW, WB, BW, BB)$ & printed d100 bands & exact die \\
\midrule
corr & $(\tfrac12, 0, 0, \tfrac12)$ & 1--50 / --- / --- / 51--100 & \\
anti & $(0, \tfrac12, \tfrac12, 0)$ & --- / 1--50 / 51--100 / --- & \\
unif & $(\tfrac14, \tfrac14, \tfrac14, \tfrac14)$ & 1--25 / 26--50 / 51--75 / 76--100 & \\
$e^{8}$ & $(\tfrac18, \tfrac38, \tfrac38, \tfrac18)$ & 1--12 / 13--50 / 51--88 / 89--100 & d8: 1 / 2--4 / 5--7 / 8 \\
$\bar e^{8}$ & $(\tfrac38, \tfrac18, \tfrac18, \tfrac38)$ & 1--38 / 39--50 / 51--62 / 63--100 & d8: 1--3 / 4 / 5 / 6--8 \\
$g$ & $(q_+, q_-, q_-, q_+)$ & 1--47 / 48--50 / 51--53 / 54--100 & \\
$\bar g$ & $(q_-, q_+, q_+, q_-)$ & 1--3 / 4--50 / 51--97 / 98--100 & \\
\bottomrule
\end{tabular}
\caption{The seven outcome patterns of the printed pair tables, with
$q_\pm = (2 \pm \sqrt 3)/8$. The rulebook's footnotes give the eight-sided-die
procedure for the eighth-family rows and call their d100 bands an approximation.}
\label{tab:patterns}
\end{table}

\begin{table}[H]
\centering
\begin{tabular}{llcccc}
\toprule
member 1 & member 2 & $\Phi^+$ & $\Psi^+$ & $\Phi^-$ & $\Psi^-$ \\
\midrule
$I$ & $I$ & corr & anti & corr & anti \\
$H$ & $I$ & unif & unif & unif & unif \\
$U$ & $I$ & $e^{8}$ & $\bar e^{8}$ & $e^{8}$ & $\bar e^{8}$ \\
$UH$ & $I$ & $g$ & $\bar g$ & $g$ & $\bar g$ \\
$I$ & $U$ & $e^{8}$ & $\bar e^{8}$ & $e^{8}$ & $\bar e^{8}$ \\
$H$ & $U$ & $g$ & $g$ & $\bar g$ & $\bar g$ \\
$U$ & $U$ & corr & $\mathbf{\bar e^{8}}$\,$^\dagger$ & $\mathbf{e^{8}}$\,$^\dagger$ & anti \\
$UH$ & $U$ & $\bar g$ & unif & unif & $g$ \\
\bottomrule
\end{tabular}
\caption{The 32 printed pair rows (8 marker configurations $\times$ 4 Bell cores),
every entry the exact Born distribution of the decorated state. All 32 printed rows
match these patterns within the rounding of Table~\ref{tab:patterns} --- except the two
marked $^\dagger$, which the rulebook prints as anti and corr respectively: erratum E2,
Section~\ref{sec:erratum}. Note the $(H, U)$ row grouping by \emph{sign} rather than by
colour type --- non-obvious, and exactly right.}
\label{tab:pairs}
\end{table}

\begin{warningbox}
\textbf{The rounding bound, audited.} Every printed d100 band is an integer, so the
exact values $\{0, 1, \tfrac12, \tfrac14, \tfrac34, \tfrac18, \tfrac38, p_\pm, q_\pm\}$
appear as $\{0, 100, 50, 25, 75, 12, 38, 93, 7, 47, 3\}$ out of 100. Over all 38
conforming rows (8 single-cell + 30 pair rows) the largest gap between print and physics
is \textbf{exactly $0.005$} --- half a percentage point, attained by the $12.5 \to 12$ and
$37.5 \to 38$ roundings; the irrational rows err by $0.0030$ (93/7) and $0.0035$ (47/3).
For the eighth-family rows the rulebook itself supplies an exact eight-sided-die
procedure and calls the d100 version an approximation --- the source documents its own
rounding layer. The digital game's exact mode extends that exactness to every row,
including the irrational ones no die can realize.
\end{warningbox}

\begin{industrybox}
\textbf{Randomness as a resource.} The measurement phase is a hand-executed Born
sampler --- the same physics that commercial quantum random-number generators package in
a chip: a superposed state, a projective measurement, an outcome no algorithm predicts.
The dice toggle is that industry's central distinction made playable: printed-tables
mode is a fixed integer approximation of the distribution (a pseudo-randomness of the
values), exact mode samples the true Born probabilities. On 38 of 40 rows you cannot
tell them apart to better than half a percentage point; on the two erratum rows only
the physics is served.
\end{industrybox}

\section{Part 5: Measurement --- The Only Move You Can't Take Back}

\begin{playbox}
Open the \textbf{Sandbox} tab. Reset the cell to $W$ and press the guided chain:
\emph{orientation $\to$ colour $\to$ orientation}. Run it many times and read the
counters: \textbf{50/50 at every step}. The white tile you started from is gone after
the first orientation measurement --- and measuring the colour again does not bring it
back. Then run both orientation procedures on the same decorated tile and compare what
is left on the cell afterwards.
\end{playbox}

\begin{questionbox}
You just measured the tile white. Measure its orientation, then look at the colour
again: half your whites turned black. Who repainted the tile?
\end{questionbox}

The endgame sweep visits cells $1 \to 9$. Monochrome tiles are skipped; every superposed
or decorated tile resolves by one roll against its printed row; a triangle pair resolves
\emph{jointly} --- one roll for both cells --- when the sweep first reaches it (at its
lower-numbered cell; the partner is monochrome by the time the sweep arrives there).
Every resolved state is replaced by a plain, undecorated outcome tile. Sweep again and
nothing changes and no dice are consumed: that is collapse, operationally --- outcomes,
once real, are stable.

\begin{shockbox}
\textbf{Measurement is the game's only irreversible move.} Every card can be undone
--- $X$, $Y$, $Z$, $H$ and the pink square are their own inverses, $C_X$ un-entangles
what it entangled. The dice cannot be un-rolled, and a white tile carries no memory of
the bianconera it came from. Reversible dynamics, irreversible measurement: the deepest
asymmetry in quantum mechanics, playable with cardboard.
\end{shockbox}

The \textbf{sandbox} lifts the rulebook's single-cell solitaire into a lab instrument.
A colour measurement asks the $\ket{0}/\ket{1}$ question. An \emph{orientation}
measurement asks the $\ket{+}/\ket{-}$ question, and the game offers it in two
equivalent-but-different ways, named as in Lab 5.

\begin{procedurebox}
\textbf{Two ways to ask the orientation question.}
\begin{itemize}
    \item \textbf{Procedure 1} projects directly onto the orientation states
          $\{\ket{+}, \ket{-}\}$: probabilities $|\braket{\pm}{\psi}|^2$, and the
          post-measurement state is an oriented bianconera tile.
    \item \textbf{Procedure 2} is the rulebook's own construction: apply $H$ --- the
          canonical rotation for this basis --- measure \emph{colour}, and relabel the
          outcomes (``swap the labels on the apparatus''). The post-measurement state is
          a monochrome tile.
\end{itemize}
Identical statistics, different post-states. The statistics agree because
$|\bra{0}H\ket{\psi}|^2 = |\braket{+}{\psi}|^2$ for every $\psi$ --- checked in the
sandbox on a state where the answer is \emph{not} 50/50: for the decorated tile
$U\ket{+} = R(-15^\circ)\ket{0}$ both procedures give $P(+) = \tfrac14$,
$P(-) = \tfrac34$ (since $\cos 15^\circ - \sin 15^\circ = \sqrt{1 - \sin 30^\circ}$).
The post-states differ: the overlap between Procedure 1's $\ket{+}$ tile and
Procedure 2's $\ket{0}$ tile is $|\braket{+}{0}|^2 = \tfrac12$, not 1. Downstream
cards see different tiles --- the Lab 5 and Lab 7 distinction, now on cardboard.
\end{procedurebox}

Colour and orientation are \emph{incompatible}: measuring one completely randomizes the
other. Starting from a white tile, the chain orientation $\to$ colour $\to$ orientation
gives 50/50 at every step. Nobody repainted the tile --- the colour question simply has
no answer while the orientation is definite, and vice versa.

\begin{numberbox}
$|\braket{0}{+}|^2 = |\braket{+}{0}|^2 = \tfrac12$ --- the two bases are mutually
unbiased (Lab 6's key-distribution pair, in tile form), so the chain reads $0.500$ at
every step, forever. In the sandbox: run the guided chain from $W$; then run both
procedures on the same decorated tile and compare post-states. And for the scoring
rule, three boards you can set up and verify: an all-white board scores 8--0 (every
line); white on the odd cells 1, 3, 5, 7, 9 and black elsewhere scores 2--0 (both
diagonals); a white top row over a black remainder scores 1--2.
\end{numberbox}

\begin{industrybox}
\textbf{Stern--Gerlach as a solitaire.} Replace colour/orientation by spin-$z$/spin-$x$
and the sandbox chain is the sequential Stern--Gerlach experiment: a spin-up beam,
measured along $x$, then along $z$ again, is 50/50 at every stage. That is exactly the
readout principle of today's spin-qubit hardware, which reads one electron's spin at a
time and pays the same price --- a measurement along one axis erases the other. The
game's single cell is the same physics with dice.
\end{industrybox}

\begin{connectionbox}
Lab 6 needed two bases that reveal nothing about each other --- $Z$ and $X$ --- so that
an eavesdropper measuring in the wrong one both learns nothing and leaves a trace. The
sandbox chain is that property on one tile: colour and orientation are exactly those two
bases, and ``half your whites turned black'' is the trace.
\end{connectionbox}

\section{The Erratum Section: When Formalism Referees the Referee}
\label{sec:erratum}

This adaptation verified every printed measurement row against the underlying quantum
mechanics. The result honours the game: \textbf{38 of the 40 printed rows are exact
within their own stated rounding} --- the largest gap is half a percentage point, and the
authors even rounded 12.5/37.5 in opposite directions (12 and 38) so the row totals and
marginals stay perfect. Three small slips remain, declared here openly, because the
game's own philosophy --- the rules \emph{are} the physics --- deserves to be taken
seriously.

\begin{shockbox}
\textbf{The printed game itself slipped --- and the formalism catches it.} The
double-pink-square table claims two pink squares never matter on an entangled pair.
That is exactly true for two of the four pair types --- and provably impossible for the
other two. Rotational invariance is a property of \emph{specific states}, not of
entanglement in general. An engine that checks every number caught a table that
playtesting could not.
\end{shockbox}

\begin{erratumbox}
\textbf{E2 --- the double-$U$ rows.} The printed table for pairs with a pink square on
\emph{each} member (the rulebook's Fig.~2.16, top) shows all four pair types unchanged
--- perfect correlation or anticorrelation, as if the squares cancelled. Here is the
derivation, step by step.

\medskip
\noindent\textbf{Step 1 --- the matrix picture.} Write a two-cell amplitude vector as a
$2 \times 2$ matrix $M$ whose entry $M_{jk}$ is the amplitude of $\ket{jk}$. Then
$\ket{\Phi^+} \leftrightarrow I/\sqrt 2$, $\ket{\Phi^-} \leftrightarrow Z/\sqrt 2$,
$\ket{\Psi^+} \leftrightarrow X/\sqrt 2$, $\ket{\Psi^-} \leftrightarrow ZX/\sqrt 2$, and
a product operation acts as
\begin{equation}
    (A \otimes B)\ket{\Phi^+} \;\leftrightarrow\; A\,B^{T}/\sqrt 2,
    \qquad\text{so}\qquad
    (U \otimes U)\ket{\Phi^-} \;\leftrightarrow\; U Z U^{T}/\sqrt 2 .
    \label{eq:matrixpic}
\end{equation}
\textbf{Step 2 --- the two rows that are right.} $U U^{T} = I$ for any rotation, so
$(U \otimes U)\ket{\Phi^+} = \ket{\Phi^+}$ exactly (Lab 7's ``rotate both together''
invariance). And $U\,(ZX)\,U^{T} = ZX$ for every rotation, because $ZX$ is itself the
rotation by $90^\circ$ and rotations commute --- so the singlet $\ket{\Psi^-}$ is
invariant too. Rows 1 and 4 of the printed table are exactly right.

\medskip
\noindent\textbf{Step 3 --- the two rows that cannot be.} For the printed $\Phi^-$
row (perfect correlation) the matrix $U Z U^{T}$ would have to be \emph{diagonal}. But for
$U = R(\theta)$,
\begin{equation}
    R(\theta)\,Z\,R(\theta)^{T} = \begin{pmatrix} \cos 2\theta & \sin 2\theta \\
    \sin 2\theta & -\cos 2\theta \end{pmatrix},
    \qquad \theta = -60^\circ:\quad
    U Z U^{T} = \begin{pmatrix} -\tfrac12 & -\tfrac{\sqrt 3}{2} \\
    -\tfrac{\sqrt 3}{2} & \tfrac12 \end{pmatrix},
    \label{eq:uzut}
\end{equation}
so the off-diagonal amplitude is $-\tfrac{\sqrt 3}{2}\cdot\tfrac{1}{\sqrt 2}$ and
\begin{equation}
    \bigl|(U Z U^{T})_{01}\bigr|^2 = \sin^2 120^\circ = \tfrac34 \;\ne\; 0 .
    \label{eq:obstruction}
\end{equation}
The exact distribution of the double-$U$ $\Phi^-$ pair is therefore
$(\tfrac18, \tfrac38, \tfrac38, \tfrac18)$ over $(WW, WB, BW, BB)$; the same computation
with $X$ in place of $Z$ gives $(\tfrac38, \tfrac18, \tfrac18, \tfrac38)$ for the
double-$U$ $\Psi^+$ pair. Against the printed $(\tfrac12, 0, 0, \tfrac12)$ and
$(0, \tfrac12, \tfrac12, 0)$ that is a deviation of exactly $\tfrac38$ in every outcome
--- not a rounding of anything.

\medskip
\noindent\textbf{Step 4 --- no other gate would help.} Suppose \emph{some} unitary
$U$ made both the printed $\Phi^+$ row and the printed $\Phi^-$ row true. Then
$U U^{T}$ and $U Z U^{T}$ would both be diagonal, hence so would their average
$U\,\tfrac{I + Z}{2}\,U^{T} = U\ketbra{0}{0}U^{T} = u\,u^{T}$ with $u = U\ket{0}$. Its
off-diagonal entry is $u_0 u_1$, so $u_0 u_1 = 0$ and $|\bra{0}U\ket{0}|^2 \in \{0, 1\}$
--- flatly contradicting the printed 25/75 row for a decorated white tile. The same
argument rules out any \emph{pair} of gates $U \otimes V$. No assignment of any rotation,
indeed any unitary, to the pink square can make the printed table true.

\medskip
\noindent\textbf{Resolution in this adaptation.} The two rows resolve by the exact
eighth-partitions on an eight-sided die --- $\Psi^+$: 1--3 / 4 / 5 / 6--8,
$\Phi^-$: 1 / 2--4 / 5--7 / 8 --- in \emph{both} dice modes (the same exact-die
procedure the rulebook itself provides for its other eighth-valued rows); the printed
bands for these two rows are shown only in the interactive page's erratum panel, as
documentation. Framing matters: this is a transcription slip in an otherwise remarkably
accurate physical game --- two rows out of forty, in the hardest table of the book. The
design is sound; the formalism simply referees everyone, authors included. That is the
point of formalism.
\end{erratumbox}

\begin{erratumbox}
\textbf{E1 --- one number in a margin note.} A collapse-stability note late in the
rulebook (p.~37) quotes probability $\tfrac14$ for a decorated right-oriented
bianconera turning up white. The book's own printed table (the 1--93 row) and the
physics give $\tfrac{2 + \sqrt 3}{4} \approx 0.933$; the quoted $\tfrac14$ belongs to the
decorated \emph{white} tile two rows up. An authorial slip with no effect on any rule
--- the note's actual argument (collapsed outcomes are stable) is untouched. Recorded so
nobody re-quotes the number.
\end{erratumbox}

\begin{erratumbox}
\textbf{E3 --- two swapped words in one caption.} One early figure caption (the
rulebook's Fig.~2.3) names its right-oriented cell ``left-oriented'' and vice versa ---
the words are swapped relative to the drawing. Every other orientation-bearing figure
and passage in the book (five drawings, four independent prose passages) agrees with the
drawing, so the drawing is authoritative and this adaptation's tiles follow it:
right-oriented means \emph{white half on the right}. Worth recording because readers
comparing this guide against that one caption would otherwise think the adaptation had
it backwards. No probability is affected.
\end{erratumbox}

\section{Declared Deviations from the Printed Game}
\label{sec:deviations}

Everything this adaptation does differently from the physical edition, in one list:

\begin{enumerate}
    \item \textbf{Ambiguity resolutions.} Where the printed rules are silent --- the
          exact mulligan procedure, tie handling, who is control and who is target when
          $C_X$ is played, how stacked markers compose, how the d100 is read, the
          three-pair limit --- this adaptation fixes each gap physics-first, guided by
          the rulebook's own examples, and enforces the result consistently. Notably: the
          player \emph{declares} control and target on every $C_X$; marker stacks
          compose with the pink square outermost (the unique order matching the printed
          47/3 rows); ties re-roll; equal scores are a draw.
    \item \textbf{Exact mode.} The physical game rolls d100 against printed integer
          bands. This adaptation adds an optional exact mode that samples the true
          probabilities in full precision --- and documents that on 38 of 40 rows the
          printed bands sit within $0.005$ of the truth. The printed-tables mode remains
          the default and reproduces the published bands verbatim.
    \item \textbf{English translation.} The original is Italian; this guide and the
          interactive page are English, with the glossary of Table~\ref{tab:glossary}.
          Rule semantics are unchanged in translation.
    \item \textbf{A seam for a computer opponent.} The engine is a pure state machine
          with no assumptions about who chooses moves, so a vs-computer mode can attach
          later. It is deliberately \emph{not} implemented: hot-seat play between two
          humans is the mode the original game is designed around.
    \item \textbf{The two erratum rows are not playable as printed.} The double-$U$
          rows for $\Psi^+$ and $\Phi^-$ resolve by the exact eighth-partitions in
          \emph{both} dice modes; the impossible printed bands appear only in the
          erratum panel. Serving the physics --- openly --- was judged more faithful to
          the game's philosophy than serving the typo.
\end{enumerate}

\noindent For completeness, the game's own numbers as implemented: the four deck
presets are Base ($I$ 6, $X$ 12, $Y$ 6, $Z$ 14, $H$ 18), +U (5, 10, 5, 11, 11, $U$ 10),
+C$_X$ (5, 10, 5, 11, 11, $C_X$ 10) and +U+C$_X$ (4, 8, 4, 9, 9, $U$ 9, $C_X$ 9); six
cards are dealt to each player, at most one mulligan each, then ten turns of draw two /
play two, so the deck is never exhausted ($12 + 20 + 12 = 44$ cards at most, counting a
worst-case mulligan, against the smallest preset); turn order and the colour White go to
the higher of two eight-sided dice; at most three entangled pairs coexist on the board.

\newpage

\section{Part 6: Guided Activities}

\begin{simulatorbox}
All activities below run on the game and sandbox of the lab's interactive page
(\texttt{lab8\_qtris.html}); the preset to use is named in each activity. All runs
are seeded --- the same seed always reproduces the same game. Activity 6 uses the
companion Qiskit notebook:\\[4pt]
\url{https://colab.research.google.com/github/LMrilo/quantum-labs-2026/blob/main/lab8/lab8_qtris.ipynb}
\end{simulatorbox}

\subsection{Activity 1: The Bianconera Tally and the \texorpdfstring{$H$}{H} Rescue (Preset \texorpdfstring{+C$_X$}{+CX})}

\begin{studentbox}
\begin{enumerate}
    \item Play several games in printed-tables mode and read the session tally: what
          fraction of half-and-half tiles came up white? Explain the number with
          equation~\eqref{eq:tiles} and the Born rule.
    \item In the operations phase, play $H$ on a right-oriented bianconera and run the
          sweep. Why does this tile consume no dice roll? Which state was it in before,
          and which after? (Equation~\eqref{eq:hclose}.)
    \item Play $Z$ on a bianconera, then $H$, then measure. Compare with $H$ alone.
          Where was the difference stored while the tile still measured 50/50?
\end{enumerate}
\end{studentbox}

\subsection{Activity 2: The Pair Counter (Preset \texorpdfstring{+C$_X$}{+CX})}

\begin{studentbox}
\begin{enumerate}
    \item Entangle two far-apart cells with $C_X$ (bianconera control, monochrome
          target). Over several games, verify that the pair counter reads 100\%
          agreement while each member's own tally stays near 50/50.
    \item Using the recipe of equation~\eqref{eq:bellrecipe}, prepare each of the four
          Bell pairs from tiles and check, with the probability preview, which two give
          $WW/BB$ and which two give $WB/BW$. Why is the sign of the pair invisible to
          the colour measurement?
    \item Play $C_X$ again on the pair you made. What comes back, and who chooses where?
          Starting from a white and a black tile, prepare the $\Psi^-$ pair using only
          cards --- one solution uses two cards.
\end{enumerate}
\end{studentbox}

\subsection{Activity 3: The 93 and Exact Mode (Preset +U)}

\begin{studentbox}
\begin{enumerate}
    \item Decorate a right-oriented bianconera with the pink square and measure it in
          printed-tables mode: name the row consulted. Switch to exact mode and measure
          again: read the probability the panel serves and compare with
          equation~\eqref{eq:cos15}.
    \item Decorate a white tile, then a black tile. Verify 25\% and 75\% against
          Table~\ref{tab:single}, and explain both from the matrix~\eqref{eq:ugate}.
    \item Why can no card stack ever make a pink-decorated tile \emph{certain}? (Its
          four possible values of $P(W)$ are $\tfrac14$, $\tfrac34$ and
          $(2 \pm \sqrt 3)/4$.) What single card removes the problem?
\end{enumerate}
\end{studentbox}

\subsection{Activity 4: Making the Erratum Visible (Preset \texorpdfstring{+U+C$_X$}{+U+CX})}

\begin{studentbox}
\begin{enumerate}
    \item Build a $\Phi^-$ pair (left-oriented control on a white target), decorate
          \emph{both} members with pink squares, and measure many times in either dice
          mode. Compare your tallies with the exact eighths
          $(\tfrac18, \tfrac38, \tfrac38, \tfrac18)$ and with the printed row shown in
          the erratum panel.
    \item Do the same with a $\Phi^+$ pair. Why do two pink squares cancel here and not
          on $\Phi^-$? Point to the line of Section~\ref{sec:erratum} that decides it.
    \item State in one sentence why \emph{no} choice of gate for the pink square could
          make the printed double-$U$ table right (Step~4 of erratum E2).
\end{enumerate}
\end{studentbox}

\subsection{Activity 5: The Sandbox Chain and the Two Procedures}

\begin{studentbox}
\begin{enumerate}
    \item Run the guided chain $W \to$ orientation $\to$ colour $\to$ orientation many
          times. Record the counters at each step and explain the three 50/50 results
          with mutual unbiasedness.
    \item On the same decorated tile $U\ket{+}$, run Procedure 1 and Procedure 2 many
          times each. Confirm the statistics agree (and are not 50/50). Then compare the
          tiles left on the cell afterwards.
    \item Try to find a state for which Procedures 1 and 2 give different
          \emph{statistics}. (You can't --- that is the theorem. Then say precisely
          what \emph{does} differ.)
\end{enumerate}
\end{studentbox}

\subsection{Activity 6: The Game in Qiskit}

The companion notebook rebuilds this guide's states as quantum circuits: the eight
single-cell tiles (the pink square as an explicit $R(-60^\circ)$, also written as a
standard $R_y$ gate), the four Bell pairs prepared exactly the way the $C_X$ card
prepares them, the full printed measurement table reproduced from state-vector
probabilities --- including a computed, not asserted, verdict on the two erratum rows
--- seeded sampling runs beside the exact numbers, and the two orientation procedures
side by side. Its tile-preparation circuit is a few lines long:

\begin{lstlisting}[language=Python]
def single_cell_circuit(name: str) -> QuantumCircuit:
    """Preparation circuit of a tile state ('W', 'B', 'R', 'L', 'UW', ...)."""
    qc = QuantumCircuit(1)
    base = name[-1]
    if base == "B":
        qc.x(0)
    elif base == "R":
        qc.h(0)
    elif base == "L":
        qc.x(0)
        qc.h(0)
    if len(name) == 2:                       # the pink square, outermost
        qc.ry(-2 * np.pi / 3, 0)
    return qc
\end{lstlisting}

\noindent The pink square is $R_y(-2\pi/3)$ --- the same matrix as
equation~\eqref{eq:ugate}, pinned entrywise in the notebook before use. All notebook
experiments are seeded and reproducible; expected values are stated with full-precision
literals and tolerances.

\begin{warningbox}
\textbf{Qiskit bit-ordering note.} Qiskit orders multi-qubit bitstrings little-endian
(rightmost bit = qubit $q_0$), while this lab series writes states big-endian (qubit 0
leftmost). Every entangled pair in QTris is a two-qubit register with qubit 0 = the
lower-numbered cell, so \emph{every} displayed pair outcome must be converted to the
leftmost-first ordering of this guide. The notebook shows exactly how --- a string
reversal --- and applies it to every histogram and state vector.
\end{warningbox}

\subsection{Activity 7: Quick Check}

\begin{reflectionbox}
\textbf{Question 1:} A right-oriented bianconera is on cell 5. Which single card makes
the upcoming colour measurement \emph{certain}, and with which outcome?
\begin{itemize}
    \item[(a)] $X$ --- certain black \qquad (b) $H$ --- certain white \qquad
               (c) $Z$ --- certain white
\end{itemize}

\textbf{Question 2:} A $\Psi^+$ pair (no markers) is measured. What is the probability
that the two cells come up the \emph{same} colour?
\begin{itemize}
    \item[(a)] $1/2$ \qquad (b) $1$ \qquad (c) $0$
\end{itemize}

\textbf{Question 3:} $P(W)$ for a pink-decorated right-oriented bianconera is
$\cos^2 15^\circ$. Which printed band encodes it, and how far is the print from the
physics?
\begin{itemize}
    \item[(a)] 1--75; off by $0.005$ \qquad (b) 1--93; off by $0.003$ \qquad
               (c) 1--93; exact
\end{itemize}

\textbf{Question 4:} Why can \emph{no} choice of gate for the pink square make the
printed double-$U$ table right?
\begin{itemize}
    \item[(a)] Because the printed $\Phi^+$ and $\Phi^-$ rows together would force
               decorated colours to stay definite, contradicting the 25/75 row.
    \item[(b)] Because two rotations can never cancel.
\end{itemize}
\end{reflectionbox}

\section{Resources}

\begin{itemize}
    \item Hamma, A.; De Simone, I.; Nazzaro, M.; Viscardi, M., \textit{QTris e
          Meccanica Quantistica}, AH Quantum Lab / Communikart, online edition for the
          NQSTI site, 2024 (Italian). --- The original game: rulebook and printable
          components; the normative source for every rule and printed table in this
          guide.
    \item NQSTI --- National Quantum Science and Technology Institute,
          \url{https://www.nqsti.it}. --- Publisher of the online edition and the
          ecosystem this adaptation belongs to.
    \item Rudolph, T., \textit{Q is for Quantum}, 2017 (Italian translation
          \textit{Quanti}, Adelphi, 2020 --- the edition the rulebook recommends). ---
          The black-and-white-ball pedagogy the game's tiles descend from.
    \item Nielsen, M. A. \& Chuang, I. L., \textit{Quantum Computation and Quantum
          Information}, Cambridge University Press, 2010. --- Pauli and Clifford
          algebra, Bell states, CNOT conventions.
    \item Seskir, Z. C. et al., ``Quantum games and interactive tools for quantum
          technologies outreach and education,'' \textit{Optical Engineering}
          \textbf{61}(8), 081809 (2022). --- Where games like QTris sit in quantum
          education.
    \item Lab 5: \textit{Quantum Measurements}, Lab 6: \textit{Quantum Key Distribution
          (BB84)} and Lab 7: \textit{Quantum Optics} (this series). --- Born rule,
          Procedures 1 and 2, mutual unbiasedness, and the pair-rotation invariance
          reused in the erratum section.
    \item Companion Jupyter notebook:
          \texttt{quantum-labs-2026/lab8/lab8\_qtris.ipynb} (Colab link in Part 6).
\end{itemize}

\end{document}
